\section{成功確率と振幅増幅}
\label{sec:hhl-success-probability}

定理4.1の仮定の下で，式\eqref{eq:hhl-before-measurement}の成功成分から
\begin{equation}
 p_{\mathrm{HHL}}
 =C^2\sum_j\frac{|\beta_j|^2}{\lambda_j^2}
 =C^2\lVert A^{-1}\ket b\rVert^2
 \label{eq:hhl-general-success-probability}
\end{equation}
を得る．固有ベクトルが正規直交するため，異なる$j$の交差項は零になる．

$A$を$1/\kappa\leq\lambda_j\leq1$となるよう規格化し，$C=1/\kappa$とする．
このとき各$j$について
\[
 \frac1{\kappa^2}\leq\frac{C^2}{\lambda_j^2}\leq1
\]
である．$\sum_j|\beta_j|^2=1$より
\begin{equation}
 \frac1{\kappa^2}\leq p_{\mathrm{HHL}}\leq1
 \label{eq:hhl-success-condition-bound}
\end{equation}
となる．実際の成功確率は$\ket b$が各固有空間に持つ重みに依存する．

HHL全体を状態準備演算$\mathcal A_{\mathrm{HHL}}$とみなし，
補助量子ビットが$\ket1$である部分を成功部分空間とすれば，
第\ref{sec:amplitude-amplification}節の振幅増幅を適用できる．
成功確率を$p_{\mathrm{HHL}}$とすると必要反復回数は
$O(1/\sqrt{p_{\mathrm{HHL}}})$であり，式\eqref{eq:hhl-success-condition-bound}の
最悪の場合には$O(\kappa)$となる．各反復にはHHL回路だけでなく
その逆回路も必要であるため，このコストを全体の計算量に含めなければならない．

